% Actual finite-II generator/reference-frame input for D.2.45.
\subsubsection*{The finite-II core three-manifold and its reference parity}

\paragraph{Parity of an actual tangent parallelization.}
For every connected closed oriented three-manifold $N$ and every
actual tangent three-frame $\tau$, its stabilized framed-bordism
class satisfies
\[
 [N,\tau]_{\rm fr}\bmod2
                  =1+\dim_{\mathbb F_2}H_1(N;\mathbb F_2).
 \tag{D.FN1}
\]
The reduction is the coefficient of a generator
$\nu\in\Omega^{\rm fr}_3=\mathbb Z/24$ modulo two.

Here is a proof retaining the chosen spin structure.  Kaplan's
Theorem~3.1 gives a spin-compatible filling $W$ with one zero-handle
and $n$ two-handles, preserving the almost framing induced by
$\tau$ \cite[Definition 1.7, Remark 1.8, and pp.243--246]{mapsKaplan79}.
The bundle $TW$ is trivial: it has a two-dimensional spine and
$w_2(W)=0$.  On its boundary use the full four-frame consisting
of the outward normal followed by $\tau$.  Its stable relative
framing obstruction is
$\lambda(TW,\tau)\in H^4(W,N;\mathbb Z)$.
Comparison with a global frame of $TW$ identifies its evaluation,
up to sign, with the stable $SO$ degree of the boundary change of
frame.  The spin structures agree, so this change lifts to
$\operatorname{Spin}$ and has no degree-one ambiguity.
The $J$-homomorphism sends that integer to the coefficient of
$[N,\tau]_{\rm fr}$ in $\mathbb Z/24$.

Reduction modulo two gives
\[
 \langle\lambda(TW,\tau),[W,N]\rangle
 =\langle w_4(TW,\tau),[W,N]\rangle
 =\langle e(TW,\tau),[W,N]\rangle
 =\chi(W)=1+n\pmod2.
 \tag{D.FN2}
\]
The Euler equality uses the prescribed outward first vector.
The integral intersection matrix of $W$ is even because $W$ is
spin.  Its reduction modulo two is therefore alternating and
has even rank.  The pair sequence identifies its cokernel over
$\mathbb F_2$ with $H_1(N;\mathbb F_2)$, since $H_1(W)=0$.
Consequently $n\equiv\dim H_1(N;\mathbb F_2)\pmod2$, proving
(D.FN1).  Changing the filling cannot alter the result; the
closed-spin difference has $\langle\lambda,[W_{\rm closed}]\rangle$
divisible by $24$.  In particular this proof applies to all spin
structures, including the exceptional Lie spin structure on
$T^3$.

\paragraph{The actual finite-II central lattice.}
Use the exact marked splitting of the good type-$II$ cap.
Its central four-manifold is an elliptic-torus bundle over a
two-torus.  The integral base lattice has generators
$\delta,\ell,g$ and the single relation
\[
                      m g=r\delta+n\ell,\qquad
                 m=6k,\quad n=me\pm1,\quad e\in\mathbb Z.
 \tag{D.FN3}
\]
The linear monodromy of $g$ on the elliptic torus is $B$, of
order six, whereas $\delta$ and $\ell$ have trivial linear
monodromy.  Put $\sigma=n-me\in\{1,-1\}$ and use the character
\[
 \chi(\delta)=1,\qquad
 \chi(\ell)=-\sigma r,\qquad \chi(g)=-\sigma re.
 \tag{D.FN4}
\]
It respects (D.FN3), since
$r-\sigma rn=-\sigma rme$.  It is primitive because its value
on $\delta$ is one.  Let $N$ be the connected fibre of the
resulting map of the central four-manifold to $S^1$.
Put $h=g+\sigma re\delta$ and $k=\ell+\sigma r\delta$.
They generate the kernel subject to $mh=nk$.  Choose integers
$a,b$ with $an+bm=1$ and set
\[
              \kappa=ah+bk,\qquad h=n\kappa,\quad k=m\kappa.
 \tag{D.FN5}
\]
The coprimality follows from $n=me+\sigma$.  Since $6\mid m$,
we have $a\equiv\sigma\pmod6$.  Thus $N$ is the actual
elliptic-torus mapping torus of an affine map with linear part
$B^a=B^\sigma$, namely $B$ or $B^{-1}$.
Translations do not alter its integral Wang sequence, and
$\det(B^\sigma-I)=1$; hence
\[
                    H_1(N;\mathbb Z)=\mathbb Z,
                    \qquad [N,\tau]_{\rm fr}=0\pmod2
 \tag{D.FN6}
\]
for every actual tangent parallelization $\tau$.

\paragraph{Why the reference stable frame is an actual tangent frame.}
Let $D$ be this reduced smooth multiple fibre.  The corrected
differential is surjective along $D$ and identifies its complex
normal line with the corrected target $L|_D$.  Its vertical
bundle is $V|_D=TD$, and the global abstract $\delta$ section is
an actual nonzero real circle direction here.
Since $\chi(\delta)=1$, this vector is transverse to
$N\subset D$.

Choose the global complex frame $(\delta,v_2)$ of $V$ and a
global frame of $L$ as in the differential argument.  Along
$N$, project the three real vectors
$J\delta,\operatorname{Re}v_2,\operatorname{Im}v_2$ to $TN$
along the transverse vector $\delta$.  They form an actual
tangent three-frame: together with $\delta$, the unprojected
vectors are a frame of $TD$.  The reference normal three-frame
is $\delta$ followed by the real and imaginary parts of a
horizontal lift of the $L$ frame.  Projection and splitting
change the full frame only by an upper triangular homotopy.
Therefore the stable tangent frame induced from the reference
$V\oplus L$ frame is the stabilization of the displayed actual
tangent frame.  Its parity is zero by (D.FN6).

This proves the reference parity for the actual finite-cap
three-manifold, rather than inferring it from its even-degree
regular three-torus cover.

\paragraph{The full-generator normal Thom comparison.}
In the original degree-two coordinates $(s,D,L')$, the three
finite Thom columns of D.2.30 for type $II$ are
\[
 K_\delta=(bn,m,0),\qquad
 K_\ell=(-br,0,m),\qquad K_g=(0,r,n).
 \tag{D.FN7}
\]
They pair with $N$ by $1,-\sigma r,-\sigma re$ respectively.
The ordinary-clutch relation is $D=-be\,s$; its bounding
chain is disjoint from this central finite-cap cycle.
After that elimination and the column change
$K_g'=K_g-eK_\ell$, the matrix and pairings are
\[
 \begin{array}{c|ccc}
       &K_\delta&K_\ell&K_g'\\ \hline
 (s,L')&(b\sigma,0)&(-br,m)&(0,\sigma)\\
 N\text{ pairing}&1&-\sigma r&0 .
 \end{array}
 \tag{D.FN8}
\]
Thus $b\,s$ bounds using $\sigma K_\delta$ and has pairing
$\sigma$.  The linking of $N$ with the torsion generator
$s$ is $\sigma/b$; in particular $N$ generates all of
$H_3(X)=\mathbb Z/b$, including when $b=20$.
This conclusion uses the full root-loop contribution.  In an
unreduced marked matrix, the four lattice-loop Thom columns
pair by $\chi(\lambda)$ and the affine root-loop column pairs
by the generally nonzero number $\chi(g)=-\sigma re$.
Setting that fifth value to zero would correspond to a
different three-manifold and can lose the odd-primary part.
